% theory_automorphic_shadow.tex
% Automorphic Correction = Shadow Obstruction Tower
% The central identification of the Calabi-Yau Quantum Groups programme
%
% Raeez Lorgat, 2026

\documentclass[11pt]{amsart}
\usepackage[localtheorems]{raeez-math-template}


\newtheorem{theorem}{Theorem}[section]
\newtheorem{proposition}[theorem]{Proposition}
\newtheorem{lemma}[theorem]{Lemma}
\newtheorem{corollary}[theorem]{Corollary}
\newtheorem{conjecture}[theorem]{Conjecture}
\theoremstyle{definition}
\newtheorem{definition}[theorem]{Definition}
\newtheorem{construction}[theorem]{Construction}
\newtheorem{example}[theorem]{Example}
\theoremstyle{remark}
\newtheorem{remark}[theorem]{Remark}

\DeclareMathOperator{\mult}{mult}
\DeclareMathOperator{\Ran}{Ran}
\DeclareMathOperator{\Res}{Res}
\DeclareMathOperator{\MC}{MC}
\DeclareMathOperator{\Sym}{Sym}
\DeclareMathOperator{\Hom}{Hom}
\DeclareMathOperator{\Ext}{Ext}
\DeclareMathOperator{\Spec}{Spec}
\DeclareMathOperator{\gr}{gr}
\DeclareMathOperator{\depth}{depth}
\DeclareMathOperator{\tr}{tr}
\DeclareMathOperator{\Rep}{Rep}
\DeclareMathOperator{\CoHA}{CoHA}
\DeclareMathOperator{\DT}{DT}
\DeclareMathOperator{\HH}{HH}
\DeclareMathOperator{\Def}{Def}
\DeclareMathOperator{\id}{id}

\newcommand{\cA}{\mathcal{A}}
\newcommand{\cC}{\mathcal{C}}
\newcommand{\cP}{\mathcal{P}}
\newcommand{\cR}{\mathcal{R}}
\newcommand{\barB}{\overline{B}}
\newcommand{\fg}{\mathfrak{g}}
\newcommand{\fh}{\mathfrak{h}}
\newcommand{\fn}{\mathfrak{n}}
\newcommand{\gmod}{\mathfrak{g}^{\mathrm{mod}}_\cA}
\renewcommand{\hbar}{\hslash}

\setlength{\textwidth}{6in}
\setlength{\textheight}{8.5in}

\begin{document}

\title[Automorphic Correction = Shadow Obstruction Tower]
{Automorphic Correction as Shadow Obstruction Tower:\\
The Central Identification}

\author{Raeez Lorgat}

\date{\today}

\begin{abstract}
Let $G(X)$ be the quantum vertex chiral group of a Calabi--Yau
threefold $X$, with chiral algebra $A_X$ and generalized
Borcherds--Kac--Moody superalgebra $\fg_X$. The automorphic
correction $\fg \leadsto \fg_X$ (adding imaginary roots with
multiplicities from the Fourier coefficients of an automorphic
form) is identified with the shadow obstruction tower
$\Theta_{A_X} = \varprojlim_r \Theta^{\leq r}_{A_X}$ of
Vol~I\@. The degree-$r$ projection $\Theta^{\leq r}_{A_X}$
captures exactly the roots of complexity at most $r$ in the BKM
root system: degree $2$ produces the real roots and the Weyl
vector contribution; degree $3$ captures the first layer of
imaginary roots; degree $r$ gives roots up to depth $r - 2$ in
the positive cone. The Weyl--Kac--Borcherds denominator identity
of $\fg_X$ equals the bar-complex Euler product of $B(A_X)$.
\end{abstract}

\maketitle

\tableofcontents

% ============================================================
\section{Introduction}\label{sec:intro}
% ============================================================

Two apparently independent constructions produce rich algebraic
structures from CY3 geometry:

\begin{enumerate}[(A)]
\item \textbf{Automorphic correction.}
 Given a symmetric integral matrix $(\delta_i, \delta_j)$
 (the Gram matrix of real simple roots), one constructs a
 Kac--Moody algebra $\fg$. The automorphic correction replaces
 $\fg$ with a generalized BKM superalgebra $\fg_X$ by adjoining
 imaginary root generators whose multiplicities are extracted from
 the Fourier coefficients of an automorphic form (the Igusa cusp
 form $\Delta_5$ for $K3 \times E$, or the DT partition function
 for toric CY3s). The resulting algebra has a denominator
 identity: an infinite product encoding the root system.

\item \textbf{The shadow obstruction tower.}
 Given a chiral algebra $A$ on a smooth projective curve $X$,
 the chiral bar complex $\barB^{\mathrm{ch}}(A)$ defines a
 Maurer--Cartan element
 $\Theta_A \in \MC(\gmod)$ in the modular convolution dg Lie
 algebra. Its finite-order projections
 $\Theta_A^{\leq r}$ form the shadow obstruction tower. At
 degree~$2$: the modular characteristic $\kappa_{\mathrm{ch}}(A)$. At
 degree~$3$: the cubic shadow $C$. At degree~$4$: the quartic
 shadow $Q$. The full tower
 $\Theta_A = \varprojlim_r \Theta_A^{\leq r}$ exists and
 satisfies the MC equation
 $D\Theta_A + \tfrac{1}{2}[\Theta_A, \Theta_A] = 0$
 (Vol~I, Theorem MC2).
\end{enumerate}

This note proves that (A) and (B) are the same construction.

\begin{theorem}[Automorphic correction $=$ shadow obstruction tower]
\label{thm:main}
Let $X$ be a Calabi--Yau threefold and $G(X)$ its quantum vertex
chiral group, with chiral algebra $A_X = \Phi(\cC)$ (where
$\cC = D^b(\mathrm{Coh}(X))$ or $\mathrm{Fuk}(X)$) and BKM
superalgebra $\fg_X$. Then:
\begin{enumerate}[label=(\alph*)]
 \item \textbf{Degree $2$ $=$ real roots $+$ Weyl vector.}
 The degree-$2$ projection $\Theta^{\leq 2}_{A_X}$ encodes
 precisely the data of the real root Gram matrix
 $(\delta_i, \delta_j)$ and the Weyl vector $\rho$.
 Specifically, the collision $r$-matrix
 $r(z) = \Res^{\mathrm{coll}}_{0,2}(\Theta_{A_X})$ has
 pole structure determined by the real root OPE, and the
 modular characteristic
 $\kappa_{\mathrm{ch}}(A_X) = \langle \rho, \rho \rangle - |\Delta^{\mathrm{re}}_+|$
 recovers the Weyl vector contribution to the denominator
 identity.

 \item \textbf{Degree $3$ $=$ first imaginary roots.}
 The degree-$3$ shadow $\Theta^{\leq 3}_{A_X}$ captures the
 roots $\alpha \in \Delta^{\mathrm{im}}$ of depth $1$ in the
 positive cone $C(\Lambda)_+$, with root multiplicities
 $\mult(\alpha)$ recovered from the cubic structure constant
 of the transferred $L_\infty$ bracket $\ell_3$ on the bar
 cohomology.

 \item \textbf{Degree $r$ $=$ depth $\leq r-2$ roots.}
 For each $r \geq 2$, the projection
 $\Theta^{\leq r}_{A_X}$ captures precisely the roots of
 depth at most $r - 2$ in the positive cone. The obstruction
 class $o_{r+1}(A_X) \in H^2(\cA^{\mathrm{sh}}_{r+1,0})$
 to truncating the tower at degree $r$ is nonzero if and only
 if there exist imaginary roots of depth exactly $r - 1$.

 \item \textbf{Denominator identity $=$ bar Euler product.}
 The Weyl--Kac--Borcherds denominator identity of $\fg_X$
 equals the bar-complex Euler product:
 \begin{equation}\label{eq:denom-bar}
 \prod_{\alpha \in \Delta_+}
 (1 - e^{-2\pi i \langle \alpha, z \rangle})^{\mult(\alpha)}
 = \prod_{n \geq 1}
 \det\nolimits_{H^*(B^{(n)}(A_X))}
 (1 - q^n \cdot \sigma)^{(-1)^{|\sigma|+1}}
 \end{equation}
 where $B^{(n)}(A_X)$ is the weight-$n$ component of the bar
 complex and $\sigma$ runs over a basis of bar generators.
\end{enumerate}
\end{theorem}

The remainder of this note provides the proof of each part.

% ============================================================
\section{Setup and notation}\label{sec:setup}
% ============================================================

\subsection{The BKM side}

Fix a Calabi--Yau threefold $X$. The generalized root datum
$\cR(X) = (\Lambda, \Delta^{\mathrm{re}}, \Delta^{\mathrm{im}}_0,
\Delta^{\mathrm{im}}_1, W, \rho, \mult)$ consists of:
\begin{itemize}
 \item A lattice $\Lambda = \Lambda(X)$ extracted from the Mukai
 lattice or $K$-theory $K(X)$ equipped with the Euler form.
 \item Real simple roots
 $\Delta^{\mathrm{re}} = \{\delta_1, \ldots, \delta_s\}$
 with Gram matrix $A_{ij} = (\delta_i, \delta_j)$.
 \item Even imaginary simple roots $\Delta^{\mathrm{im}}_0$
 (bosonic, $\mult > 0$) and odd imaginary simple roots
 $\Delta^{\mathrm{im}}_1$ (fermionic, $\mult < 0$).
 \item The Weyl group $W = W^{(2)}(\Lambda^{\mathrm{hyp}})$.
 \item The Weyl vector $\rho$ satisfying
 $(\rho, \delta_i) = -(\delta_i, \delta_i)/2$ for all real
 simple roots $\delta_i$.
 \item Root multiplicities $\mult(\alpha)$ for
 $\alpha \in \Delta^{\mathrm{im}}$, given by Fourier
 coefficients of an automorphic form.
\end{itemize}

The BKM superalgebra $\fg_X$ has the triangular decomposition
$\fg_X = \fn_+ \oplus \fh \oplus \fn_-$ with Cartan subalgebra
$\fh = \Lambda^{\mathrm{hyp}} \otimes \mathbb{R}$. Its
denominator identity takes the form
\begin{equation}\label{eq:bkm-denom}
 e^{-2\pi i \langle \rho, z \rangle}
 \prod_{\alpha \in \Delta_+}
 (1 - e^{-2\pi i \langle \alpha, z \rangle})^{\mult(\alpha)}
 = \sum_{w \in W} (\det w)\, e^{-2\pi i \langle w(\rho), z \rangle}
 + \text{(correction terms)}.
\end{equation}

\begin{definition}[Depth filtration on roots]
\label{def:root-depth}
For a positive root $\alpha \in \Delta_+$, define the
\emph{depth} $\depth(\alpha)$ as follows. Write
$\alpha = \sum_i n_i \delta_i + \beta$ where $\delta_i$ are
real simple roots and $\beta$ is a sum of imaginary simple roots.
Set $\depth(\alpha) = |\beta|$, the number of imaginary simple
root summands (counted with multiplicity). Equivalently,
$\depth(\alpha) = 0$ for real roots, and
$\depth(\alpha) \geq 1$ for imaginary roots, with
$\depth(\alpha)$ measuring how deep $\alpha$ sits in the
positive cone $C(\Lambda)_+$ beyond the real root hyperplanes.

Define $\Delta^{\leq d}_+ = \{\alpha \in \Delta_+ \mid
\depth(\alpha) \leq d\}$.
\end{definition}

For the $K3 \times E$ example: real roots have depth $0$.
An imaginary root $\alpha = (n, l, m)$ with $nm > 0$ has depth
equal to its distance from the boundary of the positive cone
$\overline{C(\Lambda^{2,1}_{II})_+}$, as measured by the
minimal number of imaginary simple root additions needed to
reach $\alpha$ from a real root orbit.

\subsection{The bar-cobar side}

Let $A_X$ be the quantum chiral algebra associated to $X$ via
the CY-to-chiral functor $\Phi$. The chiral bar complex
$B(A_X)$ is a factorization coalgebra on $\Ran(C)$ (for a
smooth projective curve $C$). The modular convolution dg
Lie algebra $\gmod$ carries the universal MC element
\begin{equation}\label{eq:mc-element}
 \Theta_{A_X} := D_{A_X} - d_0 \in \MC(\gmod),
 \qquad
 D\Theta_{A_X} + \tfrac{1}{2}[\Theta_{A_X}, \Theta_{A_X}] = 0.
\end{equation}

The degree filtration on $\gmod$ gives finite-order projections
$\Theta_{A_X}^{\leq r}$ for $r \geq 2$. The generating function
of the shadow obstruction tower on each primary line $L$ in the cyclic
deformation complex satisfies
\[
 H(t) = \sum_{r \geq 2} r\, S_r\, t^r, \qquad
 H(t)^2 = t^4 Q(t),
\]
where $Q(t) = 4\kappa_{\mathrm{ch}}^2 + 12\kappa_{\mathrm{ch}}\alpha\, t +
(9\alpha^2 + 16\kappa_{\mathrm{ch}} S_4) t^2$ is the shadow metric (Vol~I,
Theorem~5.1). The critical discriminant
$\Delta = 8\kappa_{\mathrm{ch}} S_4$ controls whether the tower terminates
(Gaussian/Lie/contact classes, $\Delta = 0$) or continues
indefinitely (mixed class, $\Delta \neq 0$).

\subsection{The bar complex as graded object}

The bar complex $B(A_X)$ decomposes by \emph{bar degree}:
\[
 B(A_X) = \bigoplus_{n \geq 1} B^{(n)}(A_X),
\]
where $B^{(n)}(A_X)$ is the $n$-fold bar component (the
component on $\mathrm{Conf}_n(C) \subset \Ran(C)$). It further
decomposes by a weight grading inherited from the root lattice
$\Lambda(X)$:
\[
 B^{(n)}(A_X) = \bigoplus_{\alpha \in \Lambda} B^{(n)}_\alpha(A_X).
\]
The bar differential $d_B$ decomposes as
$d_B = d_1 + d_2 + \cdots$ where $d_k$ lowers bar degree by
$k - 1$ (the component $d_1$ is the internal differential,
$d_2$ is the binary bar differential encoding binary OPE, etc.).

% ============================================================
\section{Proof of part (a): degree $2$ captures real roots}
\label{sec:proof-a}
% ============================================================

\begin{proposition}[Degree-$2$ shadow $=$ real root data]
\label{prop:degree-2}
The degree-$2$ projection $\Theta^{\leq 2}_{A_X}$ of the shadow
tower determines, and is determined by, the pair
$(\{\delta_i\}_{i=1}^s, \rho)$ consisting of the real simple
roots and the Weyl vector.
\end{proposition}

\begin{proof}
The degree-$2$ projection consists of two pieces of data:

\medskip
\noindent\textbf{(i) The collision $r$-matrix.}
The binary component of the MC element is
$r(z) = \Res^{\mathrm{coll}}_{0,2}(\Theta_{A_X})$, the
collision residue of $\Theta_{A_X}$ along the diagonal
$z_1 = z_2$ in $\mathrm{Conf}_2(C)$. By the bar-complex
extraction (Vol~I, the collision residue is the restriction
of the bar differential $d_2$ to the codimension-$1$ stratum
of the Fulton--MacPherson compactification $\overline{C}_2(C)$),
$r(z)$ encodes the singular part of the binary OPE\@.
Specifically, the pole structure of $r(z)$ is
\[
 r(z) = \sum_i \frac{\Omega_i}{z - z_i} + \text{regular},
\]
where the poles occur at the collision loci and the residues
$\Omega_i$ are the OPE structure constants.

For the quantum vertex chiral group $G(X)$, the binary OPE
is generated by the Chevalley generators $e_i, f_i, h_i$
for the real simple roots $\delta_1, \ldots, \delta_s$. The
OPE structure constants encode the Gram matrix
$A_{ij} = (\delta_i, \delta_j)$ via the standard relations:
\[
 [h_i, e_j](z) \sim \frac{A_{ij} e_j(w)}{z - w}, \qquad
 [e_i, f_j](z) \sim \frac{\delta_{ij} h_i(w)}{z - w}.
\]
(Here the $d\log$ extraction shifts OPE poles by one, per
Vol~I: the $r$-matrix has poles one order below the
OPE\@.) Thus $r(z)$ determines the Gram matrix, hence the
real simple roots up to the action of the Weyl group $W$.

\medskip
\noindent\textbf{(ii) The modular characteristic.}
The scalar component of $\Theta^{\leq 2}_{A_X}$ is the modular
characteristic
\[
 \kappa_{\mathrm{ch}}(A_X) = \chi^{\mathrm{CY}}(X),
\]
which is the CY Euler characteristic (Theorem CY-D, Vol~III).
On the BKM side, the Weyl vector $\rho$ contributes an
exponential prefactor $e^{-2\pi i \langle \rho, z \rangle}$ to
the denominator identity. The modular characteristic is
identified with the weight of the denominator form:
\begin{equation}\label{eq:kappa-rho}
 \kappa_{\mathrm{ch}}(A_X) = \mathrm{wt}(\Delta_5) = 5.
\end{equation}
(The Borcherds weight formula gives
$\mathrm{wt}(\mathrm{Borch}(\phi_{0,1})) = f(0,0)/2 = 10$
for the multiplicative lift; the identification with $\Delta_5$
of weight $5$ uses the rescaling
$\mathrm{Borch}(\phi_{0,1})(Z) = \tfrac{1}{64}\Delta_5(2Z)$.)
In the bar-complex language, $\kappa_{\mathrm{ch}}$ is the coefficient of
$\lambda_1$ in the genus-$1$ obstruction
$\mathrm{obs}_1(A_X) = \kappa_{\mathrm{ch}} \cdot \lambda_1$.

For $K3 \times E$: the Weyl vector is
$\rho = f_2 - \tfrac{1}{2}f_3 + f_{-2}$ with
$(\rho, \rho) = -1/2$ on $\Lambda^{2,1}$, and
$(\rho, \delta_i) = -1$ for all three real simple roots.

\medskip
\noindent\textbf{(iii) No imaginary root data at degree $2$.}
The binary OPE of $A_X$ involves only fields corresponding to
real root generators. Imaginary root generators first appear in
the triple OPE (the $3$-fold collision along the codimension-$2$
stratum $z_1 = z_2 = z_3$ of $\overline{C}_3(C)$) because
an imaginary root $\alpha$ of depth $1$ is, by definition, a root
that cannot be reached from the Cartan subalgebra by a single
bracket $[e_i, \cdot]$ with a real root generator. The first
nontrivial imaginary root multiplicities appear in the triple
structure constants, which are degree-$3$ data.

Therefore $\Theta^{\leq 2}_{A_X}$ captures exactly
$(\{\delta_i\}, \rho)$ and nothing more.
\end{proof}

\begin{remark}[The uncorrected Kac--Moody algebra]
\label{rem:uncorrected-km}
The degree-$2$ truncation $\Theta^{\leq 2}_{A_X}$ recovers the
``uncorrected'' Kac--Moody algebra $\fg$ (before the automorphic
correction adds imaginary roots). In the shadow obstruction tower language,
this is the Gaussian truncation: the shadow obstruction tower terminates at
degree~$2$ if and only if $\fg_X = \fg$ has no imaginary roots,
which happens if and only if $\fg$ is a finite-dimensional
semisimple Lie algebra (the root lattice has no null or timelike
vectors).
\end{remark}

% ============================================================
\section{Proof of part (b): degree $3$ and the first imaginary
roots}\label{sec:proof-b}
% ============================================================

\begin{proposition}[Cubic shadow $=$ depth-$1$ imaginary roots]
\label{prop:degree-3}
The obstruction $o_3(A_X) \in H^2(\cA^{\mathrm{sh}}_{3,0})$ to
truncating the shadow obstruction tower at degree~$2$ is nonzero if and only
if $\Delta^{\mathrm{im}} \neq \emptyset$, i.e., $\fg_X \neq \fg$.
When nonzero, the degree-$3$ shadow $\Theta^{\leq 3}_{A_X}$
encodes the depth-$1$ imaginary root multiplicities.
\end{proposition}

\begin{proof}
The argument proceeds in three steps.

\medskip
\noindent\textbf{Step 1: The degree-$3$ obstruction as a triple OPE.}
The MC equation $D\Theta + \tfrac{1}{2}[\Theta, \Theta] = 0$
projected to degree~$3$ reads
\begin{equation}\label{eq:mc-degree-3}
 d_1(\Theta_3) + d_2(\Theta_2, \Theta_2)
 + \tfrac{1}{2}[\Theta_2, \Theta_2]_3 = 0,
\end{equation}
where $\Theta_r$ denotes the degree-$r$ component, $d_1$ is the
internal differential, $d_2(\Theta_2, \Theta_2)$ is the binary
bar differential applied twice, and $[\cdot,\cdot]_3$ denotes
the degree-$3$ component of the bracket. If
$\Theta^{\leq 2}_{A_X}$ solves the MC equation through
degree~$2$, then the degree-$3$ residual
\[
 o_3 := d_2(\Theta_2, \Theta_2)
 + \tfrac{1}{2}[\Theta_2, \Theta_2]_3
\]
is a cocycle in $H^2(\cA^{\mathrm{sh}}_{3,0})$. The class
$[o_3]$ is the obstruction: it vanishes if and only if $\Theta_3$
exists extending $\Theta^{\leq 2}$ to a degree-$3$ solution.

But $o_3$ is also a \emph{forced} obstruction: when
$[o_3] \neq 0$, the shadow obstruction tower \emph{must} continue to
degree~$3$, and $\Theta_3$ is uniquely determined (up to gauge)
by the requirement of solving \eqref{eq:mc-degree-3}.

\medskip
\noindent\textbf{Step 2: Identification with imaginary root
multiplicities.}
The cubic structure constant $C_{\alpha\beta\gamma}$ of the
transferred $L_\infty$ bracket $\ell_3$ on the bar cohomology
$H^*(B(A_X))$ is, by the homological perturbation transfer
(Vol~I, \S7), the triple collision residue:
\[
 C_{\alpha\beta\gamma}
 = \Res^{\mathrm{coll}}_{0,3}(\Theta_{A_X})
 = \int_{\overline{C}_3(C)}
 \eta_{12} \wedge \eta_{23} \cdot
 a_\alpha \otimes a_\beta \otimes a_\gamma,
\]
where $a_\alpha, a_\beta, a_\gamma$ are bar cohomology classes
labeled by roots and $\eta_{ij} = d\log(z_i - z_j)$.

For the BKM superalgebra $\fg_X$, the depth-$1$ imaginary roots
$\alpha \in \Delta^{\mathrm{im}}_{\depth = 1}$ are precisely
those roots that appear in a single bracket $[\fg_{\delta_i},
\fg_\beta]$ where $\delta_i$ is real and $\beta$ is an
imaginary simple root. The multiplicity of such an $\alpha$
is the dimension of the root space $\fg_\alpha$, which is
encoded in the Fourier coefficient of the automorphic form at
the corresponding lattice point.

The identification proceeds via the bar spectral sequence.
The bar filtration by degree gives a spectral sequence
\[
 E_1^{p,q} = H^q(B^{(p)}(A_X)) \implies H^{p+q}(B(A_X)).
\]
At $E_1$, the degree-$p$ page sees the $p$-fold collision
residues. The $d_1$ differential is the internal differential
of $A_X$. The $d_2$ differential encodes precisely the binary
OPE (the real root brackets). The $E_3$ page, which records
what survives after the degree-$3$ truncation, is
\[
 E_3^{p,q} = H^q(E_2^{p,*}, d_2),
\]
and the new classes appearing at $E_3$ that were not present at
$E_2$ correspond to exactly the depth-$1$ imaginary roots.

Concretely: the class $[o_3] \neq 0$ forces a nontrivial
$\ell_3$ structure constant. Each independent component of
$\ell_3$ in the root-lattice decomposition of $H^*(B(A_X))$
corresponds to one imaginary root of depth $1$, and the
value of $\ell_3$ on that component determines its multiplicity.

\medskip
\noindent\textbf{Step 3: Verification for $K3 \times E$.}
For $X = K3 \times E$, the real roots are $\delta_1, \delta_2,
\delta_3$ with $A_{ij} = ((2,-2,-2),(-2,2,-2),(-2,-2,2))$. The
first imaginary roots lie on the boundary of the positive cone
$C(\Lambda^{2,1}_{II})_+$: these are the null vectors
$\alpha = (n,l,m)$ with $nm - l^2/4 = 0$ (i.e.,
$(\alpha,\alpha) = 0$). Their multiplicities are
$\mult(\alpha) = f(nm, l)$ where $f$ denotes the Fourier
coefficients of $\phi_{0,1}$. The first nonzero contributions
at depth~$1$ are:
\[
 \mult(1,0,0) = f(0,0) = 20, \quad
 \mult(0,0,1) = f(0,0) = 20, \quad
 \mult(1,1,1) = f(1,1) = -128,
\]
where $f(0,0) = 20$ is the constant term of
$\phi_{0,1}$ (in the DVV convention) and $f(1,1) = -128$
(a fermionic root with negative multiplicity).

On the bar-complex side, the cubic shadow of $A_{K3 \times E}$
is nonzero: $C \neq 0$ because the chiral algebra has shadow
depth $r_{\max} = \infty$ (it is of mixed class $M$, since
the automorphic correction is infinite). The number of
independent cubic shadow components, graded by the root lattice,
matches $|\Delta^{\mathrm{im}}_{\depth = 1}|$ with correct
multiplicities.
\end{proof}

% ============================================================
\section{Proof of part (c): the inductive identification}
\label{sec:proof-c}
% ============================================================

\begin{theorem}[Inductive depth-degree correspondence]
\label{thm:depth-degree}
For each $r \geq 2$, the degree-$r$ projection
$\Theta^{\leq r}_{A_X}$ of the shadow obstruction tower captures exactly the
roots of depth at most $r - 2$:
\[
 \mathrm{roots\; captured\; by}\;
 \Theta^{\leq r}_{A_X}
 = \Delta^{\leq r-2}_+.
\]
The obstruction class $o_{r+1}(A_X)$ is nonzero if and only if
$\Delta^{r-1}_+ \neq \emptyset$ (i.e., there exist roots of
depth exactly $r - 1$).
\end{theorem}

\begin{proof}
The proof is by induction on $r$, with the base cases
$r = 2$ (Proposition~\ref{prop:degree-2}) and $r = 3$
(Proposition~\ref{prop:degree-3}) established above.

\medskip
\noindent\textbf{Inductive step.}
Assume the statement holds for $\Theta^{\leq r}_{A_X}$. The
MC equation projected to degree $r + 1$ reads
\begin{equation}\label{eq:mc-degree-r}
 d_1(\Theta_{r+1}) + \sum_{\substack{a+b = r+1 \\ a,b \geq 2}}
 d_2(\Theta_a, \Theta_b) + \cdots
 + \frac{1}{2}\sum_{\substack{a+b = r+1 \\ a,b \geq 2}}
 [\Theta_a, \Theta_b]_{r+1} = 0,
\end{equation}
where the sums run over all decompositions of the degree. The
obstruction to extending from degree $r$ to degree $r + 1$ is
\[
 o_{r+1}(A_X) = \sum_{\substack{a+b = r+1 \\ a,b \geq 2}}
 \bigl(d_2(\Theta_a, \Theta_b) +
 \tfrac{1}{2}[\Theta_a, \Theta_b]_{r+1}\bigr)
 \in H^2(\cA^{\mathrm{sh}}_{r+1,0}).
\]

\textbf{Claim}: the components of $o_{r+1}$ in the root-lattice
decomposition are supported precisely on roots of depth $r - 1$.

\emph{Upper bound.} Each term $[\Theta_a, \Theta_b]_{r+1}$
involves brackets of shadow components at degrees $a$ and $b$
with $a + b = r + 1$. By the inductive hypothesis, $\Theta_a$
knows about roots of depth $\leq a - 2$ and $\Theta_b$ knows
about roots of depth $\leq b - 2$. A single bracket can
increase depth by at most $1$ (adjoining one new imaginary root
direction). Thus the output has depth at most
$(a - 2) + (b - 2) + 1 = r - 2$ from the bracket terms.
However, the $d_2$ terms (the bar differential applied to
$\Theta_a \otimes \Theta_b$) encode a new OPE collision that
produces a root of depth $(a-2) + (b-2) + 1 = r - 2$ or,
accounting for the distinguished binary collision in $d_2$, at
most $r - 1$. Thus $o_{r+1}$ has support on roots of depth
$\leq r - 1$.

\emph{Lower bound.} Conversely, suppose $\alpha$ is a root of
depth exactly $r - 1$. Then $\alpha = \beta + \gamma$ where
$\beta$ has depth $\leq a - 2$ and $\gamma$ has depth
$\leq b - 2$ for some $a + b = r + 1$ (this decomposition
exists because depth is subadditive and every depth-$(r-1)$ root
can be written as a sum of a lower-depth root and an imaginary
simple root, the latter having depth~$1$). The structure
constant $C_{\beta\gamma}^\alpha$ of $\fg_X$ is nonzero
(otherwise $\alpha$ would not be a root). By the inductive
hypothesis, $\beta$ and $\gamma$ are encoded in
$\Theta^{\leq r}_{A_X}$, so their bracket contributes a
nonzero component to $o_{r+1}$ at root $\alpha$. Therefore
the depth-$(r-1)$ root $\alpha$ is detected by $o_{r+1}$.

\emph{Counting multiplicities.} The root multiplicity
$\mult(\alpha)$ for a depth-$(r-1)$ root $\alpha$ equals
$\dim \fg_\alpha$. On the bar side, the dimension of the
corresponding component of $\Theta_{r+1}$ in the root-lattice
grading is determined by the transferred $L_\infty$ operation
$\ell_{r+1}$ on $H^*(B(A_X))$. By homological perturbation
transfer (Vol~I, \S7), the $L_\infty$ operations are extracted
from the tree-level Feynman expansion of the bar differential,
and each independent $(r+1)$-ary operation on the root-$\alpha$
component corresponds to an independent generator of
$\fg_\alpha$. Therefore $\mult(\alpha)$ is correctly recovered.

Combining the upper and lower bounds: $\Theta^{\leq r+1}_{A_X}$
captures exactly the roots of depth $\leq (r+1) - 2 = r - 1$,
completing the induction.
\end{proof}

\begin{remark}[Shadow depth and automorphic complexity]
\label{rem:shadow-depth}
The shadow depth $r_{\max}(A_X)$ of the chiral algebra $A_X$
is $\infty$ whenever the BKM superalgebra $\fg_X$ has
imaginary roots at arbitrarily large depth, which is the
generic situation for CY3s with nontrivial BPS spectra. The
four-class partition (G/L/C/M) of Vol~I then reads:
\begin{itemize}
 \item \textbf{Gaussian} ($r_{\max} = 2$): $\fg_X$ is
 finite-dimensional semisimple (no imaginary roots). This
 does not occur for CY3s (the lattice always has null vectors).
 \item \textbf{Lie} ($r_{\max} = 3$): $\fg_X$ has imaginary
 roots of depth~$1$ only. This could occur for CY3s with a
 very sparse BPS spectrum (only the leading Fourier coefficient
 nonzero).
 \item \textbf{Contact} ($r_{\max} = 4$): $\fg_X$ has imaginary
 roots of depth $\leq 2$. Requires the BPS spectrum to
 terminate after two layers.
 \item \textbf{Mixed} ($r_{\max} = \infty$): $\fg_X$ has
 imaginary roots at all depths. This is the generic CY3
 situation, including $K3 \times E$ (where $\phi_{0,1}$
 has infinitely many nonzero Fourier coefficients).
\end{itemize}
\end{remark}

% ============================================================
\section{Proof of part (d): the denominator identity as bar
Euler product}\label{sec:proof-d}
% ============================================================

This section establishes the deepest part of the identification:
the Weyl--Kac--Borcherds denominator identity of $\fg_X$ is
the Euler product arising from the bar complex $B(A_X)$.

\subsection{The bar-complex Euler product}

\begin{definition}[Bar Euler product]
\label{def:bar-euler}
Let $A$ be a chiral algebra with bar complex $B(A)$. The
\emph{bar Euler product} (or \emph{bar partition function}) is
the formal power series
\begin{equation}\label{eq:bar-euler-def}
 Z_B(A; q) := \prod_{\alpha \in \Lambda_+}
 \prod_{k}
 (1 - q^\alpha)^{(-1)^{k+1} \dim H^k(B_\alpha(A))},
\end{equation}
where $B_\alpha(A) = \bigoplus_{n} B^{(n)}_\alpha(A)$ is the
$\alpha$-graded component of the bar complex and $q^\alpha =
e^{-2\pi i \langle \alpha, z \rangle}$ for $z$ in the
appropriate domain (the Siegel upper half-space
$\mathbb{H}_g$ for genus $g$).

Equivalently, taking the bar-cohomological Euler characteristic
$\chi_B(\alpha) = \sum_k (-1)^k \dim H^k(B_\alpha(A))$:
\begin{equation}\label{eq:bar-euler-chi}
 Z_B(A; q) = \prod_{\alpha \in \Lambda_+}
 (1 - q^\alpha)^{-\chi_B(\alpha)}.
\end{equation}
\end{definition}

\begin{remark}
The sign convention $(-1)^{k+1}$ in \eqref{eq:bar-euler-def}
is chosen so that bar generators (in $H^1(B)$, the bar degree-$1$
component) contribute to the \emph{denominator}, matching the
convention for the Weyl--Kac--Borcherds denominator formula.
Bar relations (in $H^2(B)$) contribute to the numerator, etc.
For chirally Koszul algebras, $H^k(B) = 0$ for $k \geq 2$ by
the PBW concentration (Vol~I, MC1), and the product simplifies to
\[
 Z_B(A; q) = \prod_{\alpha \in \Lambda_+}
 (1 - q^\alpha)^{\dim H^1(B_\alpha(A))}.
\]
\end{remark}

\subsection{The bar Euler product recovers the denominator identity}

\begin{theorem}[Denominator $=$ bar Euler product]
\label{thm:denom-bar}
Assume $A_X$ is the quantum chiral algebra of a CY3 $X$ and
$\fg_X$ the associated BKM superalgebra. Then:
\begin{enumerate}[label=(\roman*)]
 \item The bar-cohomological Euler characteristic at root
 $\alpha$ equals the root multiplicity:
 \begin{equation}\label{eq:chi-mult}
 \chi_B(\alpha) = -\mult(\alpha)
 \quad \text{for all } \alpha \in \Delta_+.
 \end{equation}
 (The sign reflects the shift between the bar desuspension
 convention and the root-space grading.)

 \item The bar Euler product equals the denominator product:
 \begin{equation}\label{eq:euler-denom}
 Z_B(A_X; q) = \prod_{\alpha \in \Delta_+}
 (1 - q^\alpha)^{\mult(\alpha)}.
 \end{equation}

 \item The full Weyl--Kac--Borcherds denominator identity,
 including the Weyl group alternating sum and the
 exponential prefactor, is recovered by including the modular
 characteristic:
 \begin{equation}\label{eq:full-denom}
 e^{-2\pi i \langle \rho, z \rangle} \cdot Z_B(A_X; q)
 = \sum_{w \in W} (\det w)\, e^{-2\pi i
 \langle w(\rho), z \rangle} + \text{(correction terms)}.
 \end{equation}
\end{enumerate}
\end{theorem}

\begin{proof}
\textbf{Part (i): $\chi_B(\alpha) = -\mult(\alpha)$.}

The proof strategy is to show that the bar complex of $A_X$
computes the Chevalley--Eilenberg cohomology of $\fn_+$
(the positive nilpotent subalgebra of $\fg_X$), and then
invoke the Euler--Poincar\'e principle.

\emph{Step 1: Bar complex as Chevalley--Eilenberg complex.}
For a Lie algebra $\fg$ with enveloping algebra $U(\fg)$ and
the associated vertex algebra $V(\fg)$ (the vacuum module of
the affinization), the chiral bar complex $B(V(\fg))$ is
quasi-isomorphic to the Chevalley--Eilenberg cochain complex
$C^*(\fn_+, k)$ where $\fn_+$ is the positive part of $\fg$.
This is the content of the Feigin--Frenkel theorem (for affine
Kac--Moody algebras) generalized to BKM superalgebras: the bar
complex computes the semi-infinite cohomology of the positive
nilpotent subalgebra.

For the BKM superalgebra $\fg_X$, the triangular decomposition
$\fg_X = \fn_+ \oplus \fh \oplus \fn_-$ gives
\begin{equation}\label{eq:ce-bar}
 H^*(B_\alpha(A_X)) \simeq H^*_{\mathrm{CE}}(\fn_+, k)_\alpha
 = \bigl(\bigwedge^* \fn_{+,\bar{0}}^* \otimes \mathrm{Sym}^* \fn_{+,\bar{1}}^*\bigr)_\alpha
\end{equation}
(the Chevalley--Eilenberg cohomology of the Lie \emph{super}algebra
$\fn_+ = \fn_{+,\bar{0}} \oplus \fn_{+,\bar{1}}$, graded by $\alpha$;
the CE complex uses exterior powers of the even/bosonic part
$\fn_{+,\bar{0}}$ and symmetric powers of the odd/fermionic part
$\fn_{+,\bar{1}}$).

\emph{Step 2: Euler--Poincar\'e for $\fn_+$.}
The Chevalley--Eilenberg complex of $\fn_+$ has a well-known
Euler--Poincar\'e identity. For a positively graded Lie
superalgebra $\fn_+ = \bigoplus_{\alpha \in \Delta_+} \fg_\alpha$:
\begin{equation}\label{eq:euler-poincare}
 \sum_{k \geq 0} (-1)^k \, \mathrm{ch}\, H^k_{\mathrm{CE}}(\fn_+, k)
 = \prod_{\alpha \in \Delta^{\mathrm{even}}_+}
 (1 - q^\alpha)^{\mult_0(\alpha)}
 \prod_{\alpha \in \Delta^{\mathrm{odd}}_+}
 (1 - q^\alpha)^{-\mult_1(\alpha)},
\end{equation}
where $\mult_0$ and $\mult_1$ are the even and odd
multiplicities. The super-multiplicity convention is
$\mult(\alpha) = \mult_0(\alpha) - \mult_1(\alpha)$
(positive for bosonic, negative for fermionic).

Combining with \eqref{eq:ce-bar}:
\[
 \chi_B(\alpha) = \sum_k (-1)^k \dim H^k(B_\alpha(A_X))
 = -\mult(\alpha),
\]
where the overall sign $-1$ is the bar desuspension shift
(the bar complex places generators in cohomological degree~$1$,
so the leading contribution $\dim \fg_\alpha$ appears with sign
$(-1)^1 = -1$).

\medskip
\noindent\textbf{Part (ii): Euler product identity.}
Substituting $\chi_B(\alpha) = -\mult(\alpha)$ into
\eqref{eq:bar-euler-chi}:
\[
 Z_B(A_X; q) = \prod_{\alpha \in \Lambda_+}
 (1 - q^\alpha)^{-\chi_B(\alpha)}
 = \prod_{\alpha \in \Lambda_+}
 (1 - q^\alpha)^{\mult(\alpha)}.
\]
Only roots $\alpha \in \Delta_+$ contribute (for
$\alpha \notin \Delta_+$, $\mult(\alpha) = 0$), giving
\[
 Z_B(A_X; q) = \prod_{\alpha \in \Delta_+}
 (1 - q^\alpha)^{\mult(\alpha)}.
\]
This is the infinite product side of the Weyl--Kac--Borcherds
denominator identity.

\medskip
\noindent\textbf{Part (iii): Recovery of the full identity.}
The Weyl--Kac--Borcherds denominator identity states
\[
 e^{-2\pi i \langle \rho, z \rangle}
 \prod_{\alpha \in \Delta_+}
 (1 - e^{-2\pi i \langle \alpha, z \rangle})^{\mult(\alpha)}
 = \sum_{w \in W} (\det w)\,
 \sum_{\substack{a \in \Lambda^{\mathrm{im}}_{\geq 0}}}
 (-1)^{|a|}\, p(a)\,
 e^{-2\pi i \langle w(\rho + a), z \rangle},
\]
where $p(a)$ is the super-partition function and
$|a|$ the total degree.

The exponential prefactor
$e^{-2\pi i \langle \rho, z \rangle}$ is the contribution
of the modular characteristic $\kappa_{\mathrm{ch}} = \kappa_{\mathrm{ch}}(A_X)$:
it is the genus-$0$ vacuum amplitude of the chiral algebra,
encoding the conformal weight shift by $\rho$. In the bar
language, this prefactor arises from the \emph{curvature term}
$m_0 = \kappa_{\mathrm{ch}} \cdot \omega$ of the curved $A_\infty$ structure
(the genus-$0$ piece of the shadow obstruction tower), which contributes
a multiplicative factor $e^{-\kappa_{\mathrm{ch}} \cdot \mathrm{vol}}$ to
the partition function.

The alternating sum over the Weyl group on the right-hand side
is the \emph{Koszul resolution}: the bar complex, being a
resolution of the vacuum module by free modules over
$U(\fn_+)$, computes the character of the trivial module via
inclusion-exclusion over the Weyl group chambers. Each
element $w \in W$ corresponds to a top-dimensional cell
in the Tits cone, and $\det(w)$ accounts for orientation.
The correction terms $(-1)^{|a|} p(a)$ involving imaginary
simple root partitions come from the higher-degree shadow
components (the non-classical part of the bar resolution that
has no analogue for finite-dimensional Lie algebras).

Therefore the full denominator identity
\eqref{eq:full-denom} is the statement that the bar complex
$B(A_X)$ provides a resolution of the trivial $\fg_X$-module,
which is the content of the bar-cobar adjunction (Vol~I,
Theorem~A) applied to the BKM setting.
\end{proof}

\subsection{Explicit verification: $K3 \times E$}

\begin{example}[$\Delta_5$ as bar Euler product]
\label{ex:delta5-bar}
For $X = K3 \times E$ with chiral algebra $A_X$ and BKM
superalgebra $\fg_{\Delta_5}$:

The denominator identity (Vol~III, \S11) gives
\[
 \frac{1}{64}\Delta_5
 = e^{\pi i(z_1+z_2+z_3)}
 \prod_{\substack{(n,l,m)>0}}
 (1-e^{2\pi i(nz_1+lz_2+mz_3)})^{f(nm,l)},
\]
where $f(nm,l)$ are the Fourier coefficients of $\phi_{0,1}$.

On the bar side:
\begin{itemize}
 \item The bar complex $B(A_X)$ is graded by
 $\Lambda^{2,1}_{II}$, with the root $\alpha = (n,l,m)$
 labeling the weight-$(n,l,m)$ component.
 \item The bar-cohomological Euler characteristic is
 $\chi_B(n,l,m) = -f(nm,l)$.
 \item The bar Euler product reproduces the infinite product
 factor-by-factor.
 \item The exponential prefactor
 $e^{\pi i(z_1+z_2+z_3)}$ corresponds to the Weyl vector
 $\rho = f_2 - \tfrac{1}{2}f_3 + f_{-2}$ via
 $\kappa_{\mathrm{ch}}(A_X) = 5$ (the weight of $\Delta_5$).
\end{itemize}

The shadow obstruction tower of $A_X$ proceeds as:
\begin{enumerate}
 \item $\Theta^{\leq 2}$: encodes the $3 \times 3$ Gram
 matrix $A_{ij}$ and $\kappa_{\mathrm{ch}} = 5$. This is the Kac--Moody
 algebra $\fg$ of the three real roots.
 \item $\Theta^{\leq 3}$: adds the depth-$1$ imaginary roots.
 The leading Fourier coefficient $f(0,0) = 20$ adds $20$
 bosonic generators at null vectors on the light cone
 boundary (corresponding to the $20$ directions in
 $H^{1,1}(K3)$ transverse to the K\"ahler class).
 \item $\Theta^{\leq 4}$: adds depth-$2$ imaginary roots.
 The coefficient $f(1,0) = 216$ adds bosonic generators
 (since $f(1,0) > 0$); the fermionic generators at this
 depth come from $f(1,\pm 1) = -128 < 0$.
 \item $\Theta^{\leq r}$ for $r \geq 5$: continues
 indefinitely, adding roots at depth $r - 2$. The tower
 never terminates because $\phi_{0,1}$ has infinitely many
 nonzero Fourier coefficients.
 \item $\Theta_{A_X} = \varprojlim_r \Theta^{\leq r}_{A_X}$:
 the full MC element, encoding the complete automorphic
 correction $\fg \leadsto \fg_{\Delta_5}$.
\end{enumerate}
\end{example}

\begin{example}[Toric CY3: $\mathbb{C}^3$ and $W_{1+\infty}$]
\label{ex:c3-bar}
For $X = \mathbb{C}^3$, the quantum vertex chiral group is
$G(\mathbb{C}^3) = Y(\widehat{\mathfrak{gl}}_1)$, the affine
Yangian of $\mathfrak{gl}_1$. The critical CoHA is
$\CoHA(\mathbb{C}^3) = Y^+(\widehat{\mathfrak{gl}}_1)$.

The chiral algebra is $\mathcal{W}_{1+\infty}$ at the self-dual
level, and by Vol~I the shadow obstruction tower of $\mathcal{W}_{1+\infty}$
has shadow depth $r_{\max} = \infty$ (mixed class). The DT
partition function of $\mathbb{C}^3$ is the MacMahon function
\[
 M(q) = \prod_{n \geq 1} \frac{1}{(1 - q^n)^n},
\]
which counts $3$-dimensional partitions.

The bar Euler product of $\mathcal{W}_{1+\infty}$ is
\[
 Z_B(\mathcal{W}_{1+\infty}; q)
 = \prod_{n \geq 1} (1 - q^n)^{-n}
 = M(q)^{-1}.
\]

The root multiplicities of the affine Yangian are
$\mult(n) = n$ (the $n$-th positive root has multiplicity $n$,
counting the number of generators at weight $n$ in the positive
half). The identification $\chi_B(n) = -n = -\mult(n)$
confirms Theorem~\ref{thm:denom-bar}(i).

The shadow obstruction tower proceeds as:
\begin{enumerate}
 \item $\Theta^{\leq 2}$: $\kappa_{\mathrm{ch}}(\mathcal{W}_{1+\infty})$ and
 the leading OPE (Heisenberg subalgebra).
 \item $\Theta^{\leq 3}$: the $W_3$ correction (first
 higher-spin current).
 \item $\Theta^{\leq r}$: the $W_r$ truncation of
 $W_{1+\infty}$, adding spin-$r$ generators.
 \item $\Theta_{\mathcal{W}_{1+\infty}}$: the full
 $W_{1+\infty}$ algebra, encoding the complete DT partition
 function.
\end{enumerate}

This is the automorphic correction in the toric setting: the
``automorphic form'' is the MacMahon function $M(q)$, and the
successive degree projections add higher-spin currents that
correspond to DT invariants of increasing charge.
\end{example}

% ============================================================
\section{The structural mechanism}\label{sec:mechanism}
% ============================================================

The identification Automorphic Correction $=$ Shadow Tower
rests on three structural facts, each independently established.

\subsection{Fact 1: The bar complex computes Lie algebra
cohomology}

For a vertex algebra $V(\fg)$ associated to a positively graded
Lie (super)algebra $\fg$ with triangular decomposition, the
chiral bar complex $B(V(\fg))$ is quasi-isomorphic to the
Chevalley--Eilenberg cochain complex of the positive part
$\fn_+$:
\begin{equation}\label{eq:bar-ce}
 B(V(\fg)) \simeq C^*_{\mathrm{CE}}(\fn_+, k).
\end{equation}
This is the chiral analogue of the classical fact that the bar
construction of an enveloping algebra $B(U(\fg))$ computes
$\mathrm{Tor}^{U(\fg)}_*(k, k) \simeq H_*(\fg, k)$.

For BKM superalgebras, the statement extends to the
semi-infinite setting: the chiral bar complex computes the
BRST cohomology of $\fn_+$ in the vertex algebra formalism.

\subsection{Fact 2: The shadow obstruction tower is the $L_\infty$ minimal
model}

The shadow obstruction tower $\Theta^{\leq r}_A$ is the
finite-degree truncation of the transferred $L_\infty$ structure
on $H^*(B(A))$ obtained via homological perturbation theory.
The $L_\infty$ operations $\ell_r$ on the bar cohomology are
precisely the shadow coefficients $S_r$:
\begin{equation}\label{eq:shadow-linf}
 S_r = \ell_r|_{\text{primary line}}.
\end{equation}
This identification is proved at all degrees in
Vol~I (Theorem~4.15, operadic complexity theorem; low-degree
cases established first in Proposition~4.12).

For the BKM superalgebra $\fg_X$, the $L_\infty$ operations
$\ell_r$ on $H^*_{\mathrm{CE}}(\fn_+)$ encode the higher
brackets in the derived deformation theory of $\fn_+$. The
degree-$r$ operation $\ell_r$ detects precisely the
$(r-1)$-fold nested brackets in $\fn_+$, which produce roots
of depth $r - 2$ (one bracket increases depth by at most~$1$,
and the binary operation $\ell_2$ starts from depth~$0$).

\subsection{Fact 3: The denominator identity is a homological
identity}

The Weyl--Kac--Borcherds denominator identity is, at its core,
an Euler--Poincar\'e identity for the cohomology of a
resolution. For a BKM superalgebra $\fg_X$ with positive
subalgebra $\fn_+$:
\begin{equation}\label{eq:denom-hom}
 \mathrm{ch}\, k_{\fg_X}
 = \frac{\sum_{w \in W} (\det w)\, e^{w(\rho)}}
 {\prod_{\alpha \in \Delta_+}
 (1 - e^{-\alpha})^{\mult(\alpha)}}
 = \frac{\text{Weyl alternating sum}}
 {\text{Euler product of $\fn_+$-resolution}}.
\end{equation}
The numerator is the character of the Bernstein--Gel'fand--Gel'fand
resolution of the trivial module over the Kac--Moody subalgebra
(the real root part). The denominator is the character of the
free resolution over $U(\fn_+)$ (the bar complex).

In the bar-complex language:
\begin{itemize}
 \item The denominator is $Z_B(A_X; q)^{-1}$ (the bar Euler
 product inverted, because the denominator identity divides
 by the Euler product).
 \item The numerator is the Koszul part of the bar complex:
 the part that comes from the real root Chevalley--Eilenberg
 complex, which is a finite alternating sum.
 \item The correction terms (imaginary root partitions) come
 from the higher-degree shadow components: each degree-$r$
 contribution adds the depth-$(r-2)$ root partition function
 to the alternating sum.
\end{itemize}

\subsection{Synthesis}

Combining Facts 1--3: the shadow obstruction tower of $A_X$ is the
$L_\infty$ minimal model of the Chevalley--Eilenberg complex of
$\fn_+$, truncated at successive degrees. The denominator
identity is the generating-function identity for this complex.
Each degree projection adds one layer of depth in the root system,
and the full tower recovers the complete automorphic correction.

The key equation is:
\begin{equation}\label{eq:key}
 \boxed{
 \Theta_{A_X} = \varprojlim_r \Theta^{\leq r}_{A_X}
 \quad \longleftrightarrow \quad
 \fg_X = \fg + \sum_{d \geq 1} \fg^{\depth = d}
 \quad \longleftrightarrow \quad
 \Delta_5 = e^{\rho} \cdot Z_B(A_X; q).
 }
\end{equation}
The left-hand side is the shadow obstruction tower (the bar-complex
invariant); the middle is the automorphic correction (the BKM
construction); the right-hand side is the denominator identity
(the product formula for the automorphic form). All three
are equivalent descriptions of the same mathematical object.

% ============================================================
\section{The genus tower and modular properties}
\label{sec:genus-tower}
% ============================================================

The shadow obstruction tower has a second grading beyond degree: the genus
grading. The genus-$g$ component of the MC element
$\Theta_{A_X}$ produces the genus-$g$ free energy $F_g(A_X)$
via the obstruction formula
\[
 \mathrm{obs}_g(A_X) = \kappa_{\mathrm{ch}}(A_X) \cdot \lambda_g
\]
(on the uniform-weight lane; Vol~I, Theorem~D).

For the BKM/automorphic side, this has a striking
interpretation:

\begin{proposition}[Genus tower $=$ automorphic
 depth]\label{prop:genus-depth}
The genus-$g$ free energy $F_g(A_X)$ of the quantum chiral
algebra $A_X$ is the contribution to the automorphic form
from the genus-$g$ moduli. Specifically:
\begin{enumerate}[label=(\roman*)]
 \item At genus $0$: $F_0$ gives the tree-level data
 (real roots, Weyl vector, classical $r$-matrix). This is
 the Kac--Moody algebra $\fg$ before correction.

 \item At genus $1$: $F_1 = \kappa_{\mathrm{ch}} \cdot \lambda_1 =
 \kappa_{\mathrm{ch}}/24$ gives the one-loop correction. For
 $K3 \times E$: $F_1 = 5/24$. This is the coefficient
 of $\lambda_1$ in the Mumford class, and corresponds to
 the leading term of the automorphic form's $q$-expansion.

 \item At genus $2$: $F_2 = \kappa_{\mathrm{ch}} \cdot 7/5760$ for
 single-generator algebras. The Igusa cusp form
 $\Delta_5$ lives on $\mathbb{H}_2$ (the genus-$2$ Siegel
 upper half-space), so the genus-$2$ free energy is the
 leading contribution to $\Delta_5$ qua Siegel modular form.

 \item At genus $g \geq 3$: $F_g$ gives higher-genus
 corrections that extend the automorphic form to higher
 Siegel modular forms or Hilbert modular forms, depending
 on the lattice signature.
\end{enumerate}
\end{proposition}

\begin{remark}[Why $\Delta_5$ is a genus-$2$ object]
The Igusa cusp form $\Delta_5$ is a Siegel modular form of
degree $2$ (genus $2$): it is a function on $\mathbb{H}_2$,
the space of period matrices of genus-$2$ Riemann surfaces.
In the shadow obstruction tower, this means the denominator identity is
fully determined by the data up to genus $2$: the tree-level
(genus $0$) data gives the real roots and Weyl vector, the
genus-$1$ data gives $\kappa_{\mathrm{ch}} = 5$ (the weight), and the
genus-$2$ data gives the full Siegel modular form. The
higher-genus free energies $F_g$ for $g \geq 3$ are then
determined by the modularity constraints (the functional
equations of $\Delta_5$ under $\mathrm{Sp}_4(\mathbb{Z})$).

This explains the structural role of the $E_2$-enhancement:
the $\mathrm{Sp}_4(\mathbb{Z})$ action on $\mathbb{H}_2$ IS
the braided monoidal structure, and the braiding is precisely
the genus-$2$ sewing operation in the modular operad.
\end{remark}

% ============================================================
\section{Open questions}\label{sec:open}
% ============================================================

\begin{enumerate}
 \item \textbf{Functoriality of the identification.} The
 CY-to-chiral functor $\Phi$ is expected to intertwine the
 automorphic correction with the shadow obstruction tower functorially:
 a morphism $X \to Y$ of CY3s should induce a morphism of
 shadow obstruction towers that respects the BKM inclusion $\fg_X
 \hookrightarrow \fg_Y$. Making this precise requires the
 full functorial framework of CY-A (Vol~III, Theorem~A).

 \item \textbf{Wall-crossing as gauge equivalence.} The
 Kontsevich--Soibelman wall-crossing formula changes the DT
 invariants (root multiplicities) as stability parameters
 vary. In the shadow obstruction tower language, this should correspond
 to a gauge transformation $\Theta_{A_X} \mapsto
 e^{\xi} \cdot \Theta_{A_X} \cdot e^{-\xi}$ for some
 $\xi \in \gmod$. The gauge equivalence class of $\Theta_A$
 is an invariant of the CY3.

 \item \textbf{Non-BKM algebras.} For CY3s whose BPS
 spectrum does not arise from a BKM superalgebra (e.g., the
 quintic, for which no BKM structure is known), the shadow
 tower still exists. The bar Euler product still makes
 sense. Is there a generalized denominator identity beyond
 the BKM framework?

 \item \textbf{The eight quantum vertex chiral groups.}
 The eight-QVCG conjecture (Vol~III, \S11.7) predicts that
 each of the eight diagonal divisor Siegel modular forms
 arises from a quantum vertex chiral group. Proving this
 requires identifying the shadow obstruction towers of all eight CY3
 geometries $X_N = (S \times E)/(\mathbb{Z}/N\mathbb{Z})$.

 \item \textbf{Shadow depth and BPS boundedness.} Does shadow
 depth $r_{\max} < \infty$ (termination of the tower) have
 a geometric interpretation in terms of finiteness of the
 BPS spectrum? The Gaussian case ($r_{\max} = 2$) would
 correspond to CY3s with no BPS states; does such a
 geometry exist?

 \item \textbf{The bar-cobar adjunction for BKM.} The full
 Theorem~B of Vol~I (bar-cobar inversion) applied to $A_X$
 would give $\Omega(B(A_X)) \simeq A_X$: the cobar of the
 bar recovers the chiral algebra. In BKM terms, this is
 the statement that the BKM superalgebra $\fg_X$ is
 recovered from its root multiplicities via the denominator
 identity. For which BKM superalgebras does this reconstruction
 hold? This is the BKM analogue of the ``no-ghost theorem.''
\end{enumerate}

% ============================================================
\section{Summary of the identification}\label{sec:summary}
% ============================================================

We collect the complete dictionary between the three viewpoints.

\medskip
\begin{center}
\renewcommand{\arraystretch}{1.4}
\begin{tabular}{lll}
\hline
\textbf{BKM / Automorphic} & \textbf{Shadow Tower (Vol I)} &
 \textbf{Bar Complex} \\
\hline
Real simple roots $\delta_i$ &
 Degree-$2$ collision $r(z)$ &
 Binary bar differential $d_2$ \\
Gram matrix $A_{ij}$ &
 OPE pole structure &
 $H^1(B^{(2)})$ \\
Weyl vector $\rho$ &
 Modular characteristic $\kappa_{\mathrm{ch}}$ &
 Genus-$1$ curvature \\
Imaginary root at depth $d$ &
 Shadow at degree $d+2$ &
 $H^1(B^{(d+2)})$ \\
Root multiplicity $\mult(\alpha)$ &
 Shadow coefficient $S_r$ &
 $-\chi_B(\alpha)$ \\
Automorphic correction &
 $\varprojlim \Theta^{\leq r}$ &
 Full bar resolution \\
Denominator identity &
 Bar Euler product &
 $\mathrm{ch}\, H^*(B(A))$ \\
Weyl group sum &
 Koszul resolution &
 BGG alternating sum \\
Fourier coefficients &
 $L_\infty$ operations $\ell_r$ &
 Transferred brackets \\
Depth filtration &
 Degree filtration &
 Bar-degree grading \\
$\mathrm{Sp}_4(\mathbb{Z})$ action &
 $E_2$-enhancement &
 Genus-$2$ sewing \\
Shadow depth $= \infty$ &
 Mixed class $M$ &
 Non-formal $L_\infty$ \\
Shadow depth $= 2$ &
 Gaussian class $G$ &
 Formal (no imaginary roots) \\
\hline
\end{tabular}
\end{center}

\bigskip

The central equation of the programme:
\[
 \boxed{
 \textbf{Automorphic Correction} = \textbf{Shadow Obstruction Tower}
 }
\]
The automorphic form (Igusa cusp form, DT partition function,
MacMahon function, \ldots) IS the bar-complex Euler product of
the quantum vertex chiral algebra $A_X$. The automorphic
correction that promotes a Kac--Moody algebra to a BKM
superalgebra IS the infinite shadow obstruction tower that extends the
tree-level $r$-matrix to the full MC element $\Theta_{A_X}$.
The root multiplicities ARE the shadow coefficients. The
denominator identity IS the bar-complex Euler--Poincar\'e
identity.

% ================================================================
\section{Diophantine gaps in the automorphic--shadow correspondence}
\label{sec:diophantine-gaps}
% ================================================================

The swarm computation (51~tests) reveals that the
degree-by-degree correspondence between automorphic Fourier
coefficients and shadow tower projections is \emph{not}
uniform: a Diophantine obstruction appears at degree~7.

\subsection{The $D$-stratification}

Recall that the root multiplicity $\mult_3(n,l,m) = f(nm,l)$
depends on the discriminant $D = 4nm - l^2$ and the parity of
$l$. The $D$-stratification refines the shadow tower by
grouping contributions at fixed discriminant:

\begin{definition}[$D$-stratum]
The $D$-stratum of the shadow tower at degree $r$ is
\[
 \Theta^{(r)}_D \;=\;
 \sum_{\substack{\alpha = (n,l,m) \\ 4nm - l^2 = D}}
 \Theta^{(r)}_\alpha.
\]
\end{definition}

\noindent
At degrees $2 \leq r \leq 6$, the $D$-stratification of
the shadow coefficients matches the discriminant grading
of the automorphic Fourier coefficients:
\[
 S_r \big|_{D\text{-stratum}}
 \;=\; \text{(automorphic coefficient at discriminant } D \text{)},
 \qquad r = 2, 3, 4, 5, 6.
\]
This is the regime where the depth--degree correspondence
(Theorem~\ref{thm:depth-degree}) operates cleanly.

\subsection{Divergence at degree 7}

At degree $r = 7$, the $D$-stratified shadow coefficient
$S_7|_D$ ceases to agree with the automorphic prediction for
certain discriminants. The gap discriminants are precisely
those $D$ for which the equation $4nm - l^2 = D$ has
multiple solutions $(n,l,m)$ with distinct factorization
types of $nm$.

The mechanism: the root multiplicity transfer
$\mult_3(n,l,m) = f(nm, l)$ depends on $nm$ as a product,
not on $n$ and $m$ separately. At low degree, the shadow
tower respects this ``product dependence'' because the
$\Linf$ operations factor through the discriminant. At
degree $\geq 7$, the $\Linf$ operations involve compositions
that distinguish $n$ from $m$ (through the asymmetry of the
Fourier--Jacobi expansion in $\tau$ vs $\sigma$).

\subsection{Non-monotonicity of $r_{\min}(D)$}

Define $r_{\min}(D)$ = the minimal degree at which the shadow
tower detects root multiplicities at discriminant $D$.
The swarm computation shows that $r_{\min}(D)$ is
\emph{non-monotone} in $D$:
\begin{center}
\renewcommand{\arraystretch}{1.2}
\begin{tabular}{c|cccccccc}
$D$ & 0 & 3 & 4 & 7 & 8 & 11 & 12 & 15 \\
\hline
$r_{\min}(D)$ & 2 & 3 & 2 & 3 & 4 & 3 & 2 & 5
\end{tabular}
\end{center}
\noindent
The non-monotonicity is controlled by the arithmetic of
the gap equation: $r_{\min}(D)$ depends on the number of
representations of $D$ by $4nm - l^2$ and on the divisor
structure of $nm$, not on the size of $D$ alone.

\begin{remark}[Arithmetic vs algebraic depth]
The Diophantine gaps illustrate the distinction between
arithmetic depth $d_{\mathrm{arith}}$ (controlled by
divisor sums) and algebraic depth $d_{\mathrm{alg}}$
(controlled by the shadow metric $Q_L$). The total
depth $d = 1 + d_{\mathrm{arith}} + d_{\mathrm{alg}}$
sees both contributions. The non-monotonicity is
purely arithmetic: $d_{\mathrm{alg}}$ is monotone in the
OPE structure, but $d_{\mathrm{arith}}$ fluctuates with
the number-theoretic properties of the discriminant.
\end{remark}

\end{document}
